\clearpage
\section{Hearing the King and Receiving His Peace}
\zlabel{triptych:concise:commentary:start}
\subsection{Who becomes a subject of the King?}
Pilate asks whether Jesus is King of the Jews. Jesus first asks where the title comes from, and Pilate answers by naming those who delivered him. The governor is investigating an accusation against his authority. Jesus answers by describing his kingdom's origin: his servants do not fight to prevent his delivery. He then names the purpose of his coming and the hearing of his voice. The appointed Gospel ends at John 18:37, with this hearing, before Pilate's question about truth.

Augustine distinguishes the kingdom's place from its origin. Jesus says that it is not \emph{of} this world; he does not say that it is absent from the world. Augustine recalls Christ's prayer for disciples who remain in the world and the wheat and tares that coexist until the harvest. The kingdom is present until the end of the world, awaiting the final separation. Its heavenly origin is consistent with the present mixture of those who live within it.\footnote{Augustine, \emph{Tractates on John} 115.1--2.}

Augustine explains the kingdom's origin through regeneration. People formerly under the prince of this world become Christ's people through a new birth. He cites Colossians 1:13, the transfer into the beloved Son's kingdom, to explain this change in those who continue to live in the world. The Collect's wounded human families need healing of the persons whose sin divides them. With Augustine's account, their gathering under Christ's gentle dominion can be understood as the new life by which willing allegiance becomes possible.

The next difficulty lies in the order of Jesus' words about truth and hearing. Does someone first hear correctly and earn the right to belong to truth? Augustine answers that belonging is the prior gift from which obedient hearing follows. He distinguishes natural creaturehood from the grace at issue. All human beings exist through the Word, yet not all believe. The hearing Jesus names must therefore be more than the universal relation of creature to Creator. It is inward belief and obedience, enabled by the truth to which the hearer has been given.\footnote{Augustine, \emph{Tractates on John} 115.4.}

Chrysostom brings the same problem into view from the Epistle's opening thanksgiving. An inheritance may be offered to someone unable to enjoy it; here the Father also makes its recipients fit to share it. Rescue itself does not exhaust the generosity. Deliverance from darkness is followed by entrance into the Son's kingdom. The gift includes both freedom from captivity and a new royal relation. Augustine's account of hearing explains the willing response within that relation; Chrysostom's account of fitness explains why its recipients give thanks rather than claim an achievement.\footnote{Chrysostom, \emph{Homilies on Colossians} II, exposition of Col 1:12--15.}

The Creator has authority over every creature, including the person who resists him. Augustine's account concerns a further gift: the obedient hearing of someone who believes. Bellarmine distinguishes rightful authority from willing conversion when explaining the homage of all nations. These distinctions allow two different questions about membership. One concerns the grace by which a person begins to believe and obey; the other concerns the reconciliation of creatures whose existence has always depended on the King, even while they resist him.

This account of hearing has consequences for conduct. Allegiance expressed in obedience includes willingness to be corrected. Gratitude for grace also changes how one regards another person's conversion. Someone's testimony may help a hearer, but the response of faith depends on God's gift. The believer therefore has reason to praise and give thanks, as the Introit and Epistle do. The hearer acts, obeys and perseveres, while acknowledging that the life from which those acts proceed is received.

\subsection{How can the universal Lord receive a kingdom?}
The Mass speaks repeatedly of reception: the Lamb receives acclaim, Daniel's human figure receives dominion, and the Son asks the Father for the nations. These expressions can seem inconsistent with the Epistle's claim that everything exists through the Son. If his authority is universal, what could he still receive? Jerome and Chrysostom approach that difficulty from different texts, and their answers explain different parts of the same identity.

In Daniel 7, the human figure approaches the Ancient of Days after the vision of beastly kingdoms and judgment. Jerome explains the Son of Man through the Incarnation. The Son of God assumes human flesh; the gift of power and kingdom concerns that assumed condition. He reads the passage through Philippians 2, with its account of self-emptying and humiliation. The eternal Son truly enters the condition in which obedience and exaltation can be predicated of him. The explanation depends on the reality of the humanity in which he suffers and reigns.\footnote{Jerome, \emph{Commentary on Daniel}, VII, on 7:13--14, PL 25, cols. 670--671.}

Daniel 7:18 and 27 associate the saints with possession of this kingdom. Jerome explains their inheritance as heavenly and enduring. He rejects a fulfillment exhausted by Maccabean victory because it cannot account for the promise's universality and permanence; he also rejects an earthly millennium as the saints' final inheritance. The Son made man receives the kingdom, and the saints share its life. This account helps explain the Postcommunion's hope of lasting reign with Christ: those regenerated into his people persevere toward that inheritance.\footnote{Jerome, \emph{Commentary on Daniel}, VII, on 7:17--18,27.}

Chrysostom meets a related difficulty in Colossians 1:15. Calling Christ firstborn of creation might appear to place him first within a created series. He reads the explanatory clauses instead of isolating that title. All things, visible and invisible, are created through the Son. The source of the entire created order cannot be assigned to that order as its earliest member. The named spiritual powers reinforce the point: invisible dignity does not confer an exemption from creaturehood.\footnote{Chrysostom, \emph{Homilies on Colossians} III, on Col 1:15--17.}

Chrysostom also explains the continuing dependence of creatures. If separated from the Son's sustaining providence, they would cease to stand. An earthly ruler and an angel share this dependence, although their moral relations to Christ may differ. He then names the Son's condescension in becoming head of the Church: the one above all joins himself to those below. Its members receive their new life from the risen Son through whom they were created. Their reconciliation is accomplished by the one who sustains them.

Jerome explains received dominion through the Son's humanity; Chrysostom explains his uncreated primacy from the existence of all things through him. The Introit includes \emph{divinitatem}, divinity, among the things the slain Lamb is worthy to receive. Read with the Epistle, this is acclaim of the divine Son, not his acquisition of divine nature after death. The received Victorinus commentary connects the Lamb's victory over death with the opening of the testament and humanity's passage into life. It expounds the chapter without explaining every term of the antiphon; the ANF edition raises attribution and interpolation questions about the received text.\footnote{Received Victorinus, \emph{Commentary on the Apocalypse}, ch. 5, ANF VII, with the edition's elucidations.}

The Preface names the eternal Priest and universal King in its account of his self-offering. Augustine's treatment of Psalm 2:8 explains the Son's asking through his assumed humanity, sacrifice and intercession. In light of that account, the request for the nations concerns their salvation through the Son made man. His eternal dignity and his temporal work belong to one person. The reception of an inheritance can thus be understood alongside the creative lordship described in Colossians, without making the Son's divine power something acquired through death.\footnote{Augustine, \emph{Expositions of the Psalms}, Ps 2.7.}

\subsection{What breach does the Cross heal?}
The Collect names the dispersion of human families through sin. The Secret asks the Father that the Son offered in sacrifice give unity and peace among the nations. With the Epistle's account of reconciliation, these petitions concern a change in the relations of people estranged from God and one another. Colossians extends that reconciliation to things in heaven as well. The phrase raises a difficulty: what could holy angels need to have reconciled?

Chrysostom begins with the apparent contradiction. Heaven is the model of obedience invoked in the Lord's Prayer; war there seems inconsistent with that obedience. He resolves the difficulty by locating the breach between holy angels and sinful humanity. The angels oppose the dishonor human beings have done to their Lord. Their alienation concerns our wickedness, not a sin of their own for which they require redemption. Christ removes the estrangement by bringing humanity into their company.\footnote{Chrysostom, \emph{Homilies on Colossians} III, exposition of Col 1:20, discussion of ``things in the heavens.''}

He points to angels appearing joyfully among human beings and to humanity received in heaven. On earth the peace is twofold: human beings are reconciled with one another and with heaven. In heaven it concerns the reception of those formerly estranged from it. The angels' joy accompanies this reconciliation. Chrysostom consequently argues against relying on angels as substitutes for Christ. They minister within the peace the Son accomplishes; they cannot offer a saving way to God that bypasses him.

Augustine explains why someone can respond inwardly when creation through the Word does not by itself explain faith. Chrysostom explains the fellowship into which the reconciled are received. The person whose hearing begins in grace becomes part of a people whose life includes communion with heaven. The two accounts therefore distinguish how a person begins to belong from the shared life that belonging makes possible. Augustine's account of regeneration and Chrysostom's account of reconciliation concern different changes in the same recipient.

The Preface recounts the priestly King's offering before the assembly gives thanks with the angels. The Introit has already acclaimed the Lamb with words from the heavenly court. In this interpretation, the shared worship expresses the reconciliation Chrysostom describes. Human beings estranged from heaven enter its company through the Son's work. The assembly's thanksgiving with the angels can be understood as participation in that restored fellowship, following its praise of the sacrifice through which peace is made.

Chrysostom distinguishes the gift of reconciliation from the manner of its accomplishment. Humanity is reconciled to God; the Son himself does this, through his death and through the Cross. The one through whom his enemies exist submits to death for them. His conduct before Pilate can be understood in relation to that self-offering. His servants do not fight to prevent his delivery, and the Passion proceeds. John recounts the trial of the same Son whose death Colossians identifies as the means of peace.\footnote{Chrysostom, \emph{Homilies on Colossians} III, on Col 1:19--20; John 18:10--11,33--37.}

Colossians 1:21--23 turns from the cosmic account to former enemies and calls them to perseverance in faith. Chrysostom's heavenly reconciliation concerns holy angels and humanity, not universal pardon for every rebellious spiritual being. Nor does the Son's authority over all mean every person already enjoys saving friendship with him. The recipients have received a change in their relation to God and must live accordingly. Their perseverance is part of the passage's account of the reconciliation they have received.

The prayer for divided human families asks for a change among them as well. Their sin has wounded their relations with one another. A community receiving peace must therefore acknowledge the wrongdoing that divides it and seek the conversion of those responsible, including its own members. The Secret identifies Christ, the Son offered in sacrifice, as the giver of the unity and peace being sought. His self-offering makes their reconciliation possible, and those receiving it must respond to the gift.

\subsection{What does the homage of all nations mean?}
The Gradual selects Psalm 71:8 and 11, describing worldwide dominion and homage from rulers and peoples. These are two selected verses, not a continuous reading of verses 8--11. The intervening and following psalm describes homage and deliverance; verses 12--14 recount the king's rescue of the needy. The universal homage is thus paid to the ruler who cares for the helpless. His authority and his justice concern the same peoples.

Augustine and Bellarmine genuinely differ over the opening petition used by the Introit. Augustine begins with Christ as the true Solomon, the peacemaker whose kingdom surpasses that of the historical king. Bellarmine allows an immediate prayer by David for Solomon's just government. That prayer has an intelligible historical object before the psalm's scope and endurance lead beyond Solomon to Christ. Both affirm Christological fulfillment; they differ over how directly the opening speaks of it.\footnote{Augustine, \emph{Expositions of the Psalms}, Ps 71.1--2; Bellarmine, \emph{Commentary on Psalms}, Ps 71, opening and comments 4--7.}

Bellarmine's account gives the charge to govern justly an immediate place within Israel's royal life. Solomon can receive a responsibility that the prayer asks him to discharge well. The promises of endurance and universal dominion then exceed his reign and find their fulfillment in Christ. Augustine begins with that fulfillment in the true Solomon. This disagreement concerns the opening petition; the further question of whether all rulers already obey willingly requires its own answer.

Bellarmine addresses precisely that difficulty at biblical verse 11, treated in his numbered comment 10. He distinguishes three possible senses. Scriptural generalization can describe the widespread conversion of rulers and nations while allowing exceptions. The complete recognition may belong to final judgment. Or the words may concern Christ's present authority, which holds even over those who do not consent. These answers differ in subject and time: conversion concerns a willing change; judgment concerns unavoidable acknowledgment; authority concerns the King's actual right.\footnote{Bellarmine, \emph{Commentary on Psalms}, Ps 71, comment 10.}

For Augustine, the hearer believes and obeys through grace. For Chrysostom, even a creature estranged from the Son depends on him for existence. Bellarmine's alternatives distinguish that existing authority from a willing change in the subject and from the acknowledgment of judgment. Resistance therefore does not disprove Christ's authority. Yet acknowledgment compelled at judgment is different from the saving communion of those converted to him. The universal scope of his rule leaves the need for conversion intact.

Bellarmine interprets the psalm's rescue of the poor through liberation from the devil's captivity, forgiveness and restored grace. Humanity's poverty becomes the condition in which people recognize their need of the King, who does not despise them. This deliverance can be read beside Colossians' transfer from darkness: those held captive are freed and restored to life with God. Pius XI's feast teaching extends the moral claims of Christ's kingship to public life. Those who exercise authority remain responsible to the King whose justice includes care for the helpless.\footnote{Bellarmine, \emph{Commentary on Psalms}, Ps 71, comments 11--13; Pius XI, \emph{Quas primas}, nos. 12--20.}

The Offertory's Psalm 2:8 reaches the same peoples through the Son's petition. Augustine emphasizes incarnate intercession for their salvation; Bellarmine relates their inheritance to the Son's rightful possession. The following rod-of-iron verse belongs to the psalm's context, not the appointed Offertory. Bellarmine's treatment includes correction of sinful habits and final accountability. The royal image cannot simply become a commission for worshippers to coerce conversion. The Gospel's King gives himself, and his subjects are called to hear.\footnote{Augustine, \emph{Expositions of the Psalms}, Ps 2.7; Bellarmine, \emph{Commentary on Psalms}, Ps 2, comments on vv. 7--10.}

\subsection{How does received peace become a way of life?}
The Communion selects clauses from both Psalm 28:10 and 11: the Lord's eternal kingship and his blessing of peace. The preceding psalm describes his voice amid a storm. Augustine reads the conclusion as the Church's preservation amid the world's flood and the Lord's future reign in his saints. In his account, peace in Christ is received during continuing adversity. The blessing is consistent with the conflict through which the Church perseveres toward lasting life with him.\footnote{Augustine, \emph{Expositions of the Psalms}, Ps 28.10--11.}

Bellarmine preserves natural and ecclesial readings of the psalm and does not settle all alternatives concerning the flood. He distinguishes present strength from final blessing. That distinction helps explain the service and hope of those praying the Postcommunion. Already nourished by the sacrament and glorying in service under Christ's standard, they ask to reign with him in heaven. Their service can be understood as perseverance sustained by the nourishment received. Read beside the Gospel, its military image concerns fidelity to the King whose servants do not fight to prevent his Passion.\footnote{Bellarmine, \emph{Commentary on Psalms}, Ps 28, opening and conclusion.}

Chrysostom approaches the same blessing from the relation between the recipients. His exposition of heavenly reconciliation turns toward the Church's repeated exchange of peace. A greeting addressed to all cannot be made true by peace with one neighbor while a chosen feud is preserved with another. He compares that partial concord to a body with some members healthy and others disordered. The health of the whole cannot be established by pointing only to its healthy part.\footnote{Chrysostom, \emph{Homilies on Colossians} III, practical application following the exposition of heavenly reconciliation.}

Chrysostom's application follows from the reconciliation he has described. The Son came to reconcile heaven and earth, and a recipient of the Church's peace refuses reconciliation with a fellow member. For Chrysostom, deliberately preserving the feud contradicts the gift exchanged. The Secret's request for the nations and the Communion's blessing of the people can be understood beside this obligation. Those who receive the King's peace must seek its realization among the people with whom they worship.

The two interpretations lead to different moral consequences. Augustine's hearer owes the capacity to respond to grace, so patient testimony and correction of one's own habits follow from the gift rather than earn it. Chrysostom's recipient of peace must seek reconciliation with a fellow member of the Church. Gratitude concerns the beginning of the believer's life; concord concerns its exercise among others. The believer who receives truth also has to learn what reconciliation requires.

Jerome's exposition of Daniel concerns the saints' enduring heavenly inheritance. Bellarmine's final acknowledgment concerns the complete homage of every power at judgment. The Postcommunion asks for shared reign with Christ, the lasting life toward which his people persevere. That hope differs from the acknowledgment compelled of those who resist him. Saving communion and unavoidable homage remain distinct even when both acknowledge the same universal King.
